% Options for packages loaded elsewhere
\PassOptionsToPackage{unicode}{hyperref}
\PassOptionsToPackage{hyphens}{url}
\documentclass[
  spanish,
]{article}
\usepackage{xcolor}
\usepackage[margin=1in]{geometry}
\usepackage{amsmath,amssymb}
\setcounter{secnumdepth}{5}
\usepackage{iftex}
\ifPDFTeX
  \usepackage[T1]{fontenc}
  \usepackage[utf8]{inputenc}
  \usepackage{textcomp} % provide euro and other symbols
\else % if luatex or xetex
  \usepackage{unicode-math} % this also loads fontspec
  \defaultfontfeatures{Scale=MatchLowercase}
  \defaultfontfeatures[\rmfamily]{Ligatures=TeX,Scale=1}
\fi
\usepackage{lmodern}
\ifPDFTeX\else
  % xetex/luatex font selection
\fi
% Use upquote if available, for straight quotes in verbatim environments
\IfFileExists{upquote.sty}{\usepackage{upquote}}{}
\IfFileExists{microtype.sty}{% use microtype if available
  \usepackage[]{microtype}
  \UseMicrotypeSet[protrusion]{basicmath} % disable protrusion for tt fonts
}{}
\makeatletter
\@ifundefined{KOMAClassName}{% if non-KOMA class
  \IfFileExists{parskip.sty}{%
    \usepackage{parskip}
  }{% else
    \setlength{\parindent}{0pt}
    \setlength{\parskip}{6pt plus 2pt minus 1pt}}
}{% if KOMA class
  \KOMAoptions{parskip=half}}
\makeatother
\usepackage{color}
\usepackage{fancyvrb}
\newcommand{\VerbBar}{|}
\newcommand{\VERB}{\Verb[commandchars=\\\{\}]}
\DefineVerbatimEnvironment{Highlighting}{Verbatim}{commandchars=\\\{\}}
% Add ',fontsize=\small' for more characters per line
\usepackage{framed}
\definecolor{shadecolor}{RGB}{248,248,248}
\newenvironment{Shaded}{\begin{snugshade}}{\end{snugshade}}
\newcommand{\AlertTok}[1]{\textcolor[rgb]{0.94,0.16,0.16}{#1}}
\newcommand{\AnnotationTok}[1]{\textcolor[rgb]{0.56,0.35,0.01}{\textbf{\textit{#1}}}}
\newcommand{\AttributeTok}[1]{\textcolor[rgb]{0.13,0.29,0.53}{#1}}
\newcommand{\BaseNTok}[1]{\textcolor[rgb]{0.00,0.00,0.81}{#1}}
\newcommand{\BuiltInTok}[1]{#1}
\newcommand{\CharTok}[1]{\textcolor[rgb]{0.31,0.60,0.02}{#1}}
\newcommand{\CommentTok}[1]{\textcolor[rgb]{0.56,0.35,0.01}{\textit{#1}}}
\newcommand{\CommentVarTok}[1]{\textcolor[rgb]{0.56,0.35,0.01}{\textbf{\textit{#1}}}}
\newcommand{\ConstantTok}[1]{\textcolor[rgb]{0.56,0.35,0.01}{#1}}
\newcommand{\ControlFlowTok}[1]{\textcolor[rgb]{0.13,0.29,0.53}{\textbf{#1}}}
\newcommand{\DataTypeTok}[1]{\textcolor[rgb]{0.13,0.29,0.53}{#1}}
\newcommand{\DecValTok}[1]{\textcolor[rgb]{0.00,0.00,0.81}{#1}}
\newcommand{\DocumentationTok}[1]{\textcolor[rgb]{0.56,0.35,0.01}{\textbf{\textit{#1}}}}
\newcommand{\ErrorTok}[1]{\textcolor[rgb]{0.64,0.00,0.00}{\textbf{#1}}}
\newcommand{\ExtensionTok}[1]{#1}
\newcommand{\FloatTok}[1]{\textcolor[rgb]{0.00,0.00,0.81}{#1}}
\newcommand{\FunctionTok}[1]{\textcolor[rgb]{0.13,0.29,0.53}{\textbf{#1}}}
\newcommand{\ImportTok}[1]{#1}
\newcommand{\InformationTok}[1]{\textcolor[rgb]{0.56,0.35,0.01}{\textbf{\textit{#1}}}}
\newcommand{\KeywordTok}[1]{\textcolor[rgb]{0.13,0.29,0.53}{\textbf{#1}}}
\newcommand{\NormalTok}[1]{#1}
\newcommand{\OperatorTok}[1]{\textcolor[rgb]{0.81,0.36,0.00}{\textbf{#1}}}
\newcommand{\OtherTok}[1]{\textcolor[rgb]{0.56,0.35,0.01}{#1}}
\newcommand{\PreprocessorTok}[1]{\textcolor[rgb]{0.56,0.35,0.01}{\textit{#1}}}
\newcommand{\RegionMarkerTok}[1]{#1}
\newcommand{\SpecialCharTok}[1]{\textcolor[rgb]{0.81,0.36,0.00}{\textbf{#1}}}
\newcommand{\SpecialStringTok}[1]{\textcolor[rgb]{0.31,0.60,0.02}{#1}}
\newcommand{\StringTok}[1]{\textcolor[rgb]{0.31,0.60,0.02}{#1}}
\newcommand{\VariableTok}[1]{\textcolor[rgb]{0.00,0.00,0.00}{#1}}
\newcommand{\VerbatimStringTok}[1]{\textcolor[rgb]{0.31,0.60,0.02}{#1}}
\newcommand{\WarningTok}[1]{\textcolor[rgb]{0.56,0.35,0.01}{\textbf{\textit{#1}}}}
\usepackage{graphicx}
\makeatletter
\newsavebox\pandoc@box
\newcommand*\pandocbounded[1]{% scales image to fit in text height/width
  \sbox\pandoc@box{#1}%
  \Gscale@div\@tempa{\textheight}{\dimexpr\ht\pandoc@box+\dp\pandoc@box\relax}%
  \Gscale@div\@tempb{\linewidth}{\wd\pandoc@box}%
  \ifdim\@tempb\p@<\@tempa\p@\let\@tempa\@tempb\fi% select the smaller of both
  \ifdim\@tempa\p@<\p@\scalebox{\@tempa}{\usebox\pandoc@box}%
  \else\usebox{\pandoc@box}%
  \fi%
}
% Set default figure placement to htbp
\def\fps@figure{htbp}
\makeatother
\ifLuaTeX
\usepackage[bidi=basic]{babel}
\else
\usepackage[bidi=default]{babel}
\fi
% get rid of language-specific shorthands (see #6817):
\let\LanguageShortHands\languageshorthands
\def\languageshorthands#1{}
\setlength{\emergencystretch}{3em} % prevent overfull lines
\providecommand{\tightlist}{%
  \setlength{\itemsep}{0pt}\setlength{\parskip}{0pt}}
\usepackage{booktabs}
\usepackage{longtable}
\usepackage{array}
\usepackage{multirow}
\usepackage{float}
\usepackage{colortbl}
\usepackage{xcolor}
\usepackage{pdflscape}
\usepackage{geometry}
\geometry{margin=2.5cm}
\usepackage{fancyhdr}
\pagestyle{fancy}
Factor Sensorial}
\fancyhead[R]{\small Tesis --- Experiencia Gastronómica}
\fancyfoot[C]{\thepage}
\renewcommand{\headrulewidth}{0.4pt}
\usepackage{caption}
\captionsetup{font=small, labelfont=bf}
\usepackage{indentfirst}
\setlength{\parindent}{1.25cm}
\usepackage{booktabs}
\usepackage{longtable}
\usepackage{array}
\usepackage{multirow}
\usepackage{wrapfig}
\usepackage{float}
\usepackage{colortbl}
\usepackage{pdflscape}
\usepackage{tabu}
\usepackage{threeparttable}
\usepackage{threeparttablex}
\usepackage[normalem]{ulem}
\usepackage{makecell}
\usepackage{xcolor}
\usepackage{bookmark}
\IfFileExists{xurl.sty}{\usepackage{xurl}}{} % add URL line breaks if available
\urlstyle{same}
\hypersetup{
  pdftitle={Objetivo 4: Factor Sensorial Más Influyente},
  pdfauthor={Martín Vargas Estrada},
  pdflang={es},
  hidelinks,
  pdfcreator={LaTeX via pandoc}}

\title{Objetivo 4: Factor Sensorial Más Influyente}
\usepackage{etoolbox}
\makeatletter
\providecommand{\subtitle}[1]{% add subtitle to \maketitle
  \apptocmd{\@title}{\par {\large #1 \par}}{}{}
}
\makeatother
\subtitle{Impacto de los factores de la experiencia gastronómica en la
toma de decisiones del comensal}
\author{Martín Vargas Estrada}
\date{30 de March de 2026}

\begin{document}
\maketitle

{
\setcounter{tocdepth}{3}
\tableofcontents
}
\section{Carga de datos y
librerías}\label{carga-de-datos-y-libreruxedas}

\begin{Shaded}
\begin{Highlighting}[]
\FunctionTok{library}\NormalTok{(dplyr)}
\FunctionTok{library}\NormalTok{(tidyr)}
\FunctionTok{library}\NormalTok{(ggplot2)}
\FunctionTok{library}\NormalTok{(knitr)}
\FunctionTok{library}\NormalTok{(kableExtra)}
\FunctionTok{library}\NormalTok{(gridExtra)}
\end{Highlighting}
\end{Shaded}

\begin{Shaded}
\begin{Highlighting}[]
\FunctionTok{load}\NormalTok{(}\StringTok{"datos\_tesis\_preparados.RData"}\NormalTok{)}
\end{Highlighting}
\end{Shaded}

El objetivo específico 4 busca identificar cuál de los tres factores
sensoriales del producto (sabor (P1), presentación (P2) y frescura (P3))
es el más influyente en la toma de decisiones del comensal, medida a
través de tres variables dependientes: satisfacción (I1, Likert 1 a 5),
intención de revisitar (I2, Likert 1 a 5) e intención de recomendar (I3,
binaria Sí/No).

Se analizan las 9 combinaciones (3 ítems \(\times\) 3 VD) mediante
correlaciones de Spearman, tablas de contingencia y visualizaciones.

\newpage

\section{Correlaciones de Spearman: P1, P2, P3 vs.~variables
dependientes}\label{correlaciones-de-spearman-p1-p2-p3-vs.-variables-dependientes}

Se calcula el coeficiente de correlación de Spearman (\(\rho_s\)) entre
cada ítem del producto y cada variable dependiente. Al tratarse de
variables ordinales, Spearman es la medida de asociación más apropiada.

\begin{Shaded}
\begin{Highlighting}[]
\NormalTok{items\_prod }\OtherTok{\textless{}{-}} \FunctionTok{c}\NormalTok{(}\StringTok{"P1"}\NormalTok{, }\StringTok{"P2"}\NormalTok{, }\StringTok{"P3"}\NormalTok{)}
\NormalTok{etiq\_prod  }\OtherTok{\textless{}{-}} \FunctionTok{c}\NormalTok{(}\StringTok{"P1"} \OtherTok{=} \StringTok{"P1 (Sabor)"}\NormalTok{, }\StringTok{"P2"} \OtherTok{=} \StringTok{"P2 (Presentación)"}\NormalTok{,}
                 \StringTok{"P3"} \OtherTok{=} \StringTok{"P3 (Frescura)"}\NormalTok{)}

\NormalTok{vars\_dep   }\OtherTok{\textless{}{-}} \FunctionTok{c}\NormalTok{(}\StringTok{"idx\_satisfaccion"}\NormalTok{, }\StringTok{"idx\_revisitar"}\NormalTok{, }\StringTok{"I3\_bin"}\NormalTok{)}
\NormalTok{etiq\_dep   }\OtherTok{\textless{}{-}} \FunctionTok{c}\NormalTok{(}\StringTok{"idx\_satisfaccion"} \OtherTok{=} \StringTok{"Satisfacción"}\NormalTok{,}
                 \StringTok{"idx\_revisitar"}    \OtherTok{=} \StringTok{"Revisitar"}\NormalTok{,}
                 \StringTok{"I3\_bin"}           \OtherTok{=} \StringTok{"Recomendar"}\NormalTok{)}

\NormalTok{corr\_resultados }\OtherTok{\textless{}{-}} \FunctionTok{data.frame}\NormalTok{()}

\ControlFlowTok{for}\NormalTok{ (item }\ControlFlowTok{in}\NormalTok{ items\_prod) \{}
  \ControlFlowTok{for}\NormalTok{ (vd }\ControlFlowTok{in}\NormalTok{ vars\_dep) \{}
\NormalTok{    test\_s }\OtherTok{\textless{}{-}} \FunctionTok{cor.test}\NormalTok{(datos[[item]], datos[[vd]],}
                        \AttributeTok{method =} \StringTok{"spearman"}\NormalTok{, }\AttributeTok{exact =} \ConstantTok{FALSE}\NormalTok{)}
\NormalTok{    corr\_resultados }\OtherTok{\textless{}{-}} \FunctionTok{rbind}\NormalTok{(corr\_resultados, }\FunctionTok{data.frame}\NormalTok{(}
      \AttributeTok{Item     =}\NormalTok{ etiq\_prod[item],}
      \AttributeTok{VD       =}\NormalTok{ etiq\_dep[vd],}
      \AttributeTok{rho      =} \FunctionTok{round}\NormalTok{(test\_s}\SpecialCharTok{$}\NormalTok{estimate, }\DecValTok{3}\NormalTok{),}
      \AttributeTok{p        =}\NormalTok{ test\_s}\SpecialCharTok{$}\NormalTok{p.value,}
      \AttributeTok{Sig      =} \FunctionTok{ifelse}\NormalTok{(test\_s}\SpecialCharTok{$}\NormalTok{p.value }\SpecialCharTok{\textless{}} \FloatTok{0.05}\NormalTok{, }\StringTok{"Sí"}\NormalTok{, }\StringTok{"No"}\NormalTok{)}
\NormalTok{    ))}
\NormalTok{  \}}
\NormalTok{\}}
\end{Highlighting}
\end{Shaded}

\begin{table}[!h]
\centering
\begin{tabular}{llcc>{}c}
\toprule
Ítem Producto & Variable dependiente & $\rho_s$ & $p$ & Significativo\\
\midrule
\addlinespace[0.3em]
\multicolumn{5}{l}{\textbf{P1 (Sabor)}}\\
\hspace{1em}P1 (Sabor) & Satisfacción & 0.586 & < .001 & \textcolor{teal}{Sí}\\
\hspace{1em}P1 (Sabor) & Revisitar & 0.550 & < .001 & \textcolor{teal}{Sí}\\
\hspace{1em}P1 (Sabor) & Recomendar & 0.484 & < .001 & \textcolor{teal}{Sí}\\
\addlinespace[0.3em]
\multicolumn{5}{l}{\textbf{P2 (Presentación)}}\\
\hspace{1em}P2 (Presentación) & Satisfacción & 0.490 & < .001 & \textcolor{teal}{Sí}\\
\hspace{1em}P2 (Presentación) & Revisitar & 0.517 & < .001 & \textcolor{teal}{Sí}\\
\hspace{1em}P2 (Presentación) & Recomendar & 0.427 & < .001 & \textcolor{teal}{Sí}\\
\addlinespace[0.3em]
\multicolumn{5}{l}{\textbf{P3 (Frescura)}}\\
\hspace{1em}P3 (Frescura) & Satisfacción & 0.414 & < .001 & \textcolor{teal}{Sí}\\
\hspace{1em}P3 (Frescura) & Revisitar & 0.351 & < .001 & \textcolor{teal}{Sí}\\
\hspace{1em}P3 (Frescura) & Recomendar & 0.364 & < .001 & \textcolor{teal}{Sí}\\
\bottomrule
\end{tabular}
\end{table}

El ítem con mayor \(|\rho_s|\) promedio a través de las tres VD es el
factor sensorial más influyente.

\begin{Shaded}
\begin{Highlighting}[]
\NormalTok{resumen\_rho }\OtherTok{\textless{}{-}}\NormalTok{ corr\_resultados }\SpecialCharTok{\%\textgreater{}\%}
  \FunctionTok{group\_by}\NormalTok{(Item) }\SpecialCharTok{\%\textgreater{}\%}
  \FunctionTok{summarise}\NormalTok{(}
    \AttributeTok{rho\_medio =} \FunctionTok{round}\NormalTok{(}\FunctionTok{mean}\NormalTok{(}\FunctionTok{abs}\NormalTok{(rho)), }\DecValTok{3}\NormalTok{),}
    \AttributeTok{.groups =} \StringTok{"drop"}
\NormalTok{  ) }\SpecialCharTok{\%\textgreater{}\%}
  \FunctionTok{arrange}\NormalTok{(}\FunctionTok{desc}\NormalTok{(rho\_medio))}
\end{Highlighting}
\end{Shaded}

\begin{table}[!h]
\centering
\begin{tabular}{lc}
\toprule
Ítem Producto & $|\rho_s|$ promedio\\
\midrule
P1 (Sabor) & 0.540\\
P2 (Presentación) & 0.478\\
P3 (Frescura) & 0.376\\
\bottomrule
\end{tabular}
\end{table}

\newpage

\section{Tablas de contingencia}\label{tablas-de-contingencia}

Se presentan las tablas de contingencia cruzando cada ítem del producto
(P1, P2, P3) con cada variable dependiente. Las frecuencias muestran
cómo se distribuyen las respuestas de las VD en cada nivel del ítem.

\subsection{P1 (Sabor) vs.~Variables
dependientes}\label{p1-sabor-vs.-variables-dependientes}

\subsubsection{P1 vs.~Satisfacción
(I1)}\label{p1-vs.-satisfacciuxf3n-i1}

\begin{table}[!h]
\centering
\caption{\label{tab:cont-p1-satisf}Tabla de contingencia: P1 (Sabor) vs. Satisfacción (I1)}
\centering
\begin{tabular}[t]{lrrrrrr}
\toprule
  & 1 & 2 & 3 & 4 & 5 & Sum\\
\midrule
1 & 3 & 6 & 1 & 0 & 0 & 10\\
2 & 2 & 21 & 12 & 4 & 0 & 39\\
3 & 1 & 7 & 25 & 37 & 0 & 70\\
4 & 0 & 3 & 7 & 29 & 8 & 47\\
5 & 0 & 2 & 2 & 4 & 2 & 10\\
\addlinespace
Sum & 6 & 39 & 47 & 74 & 10 & 176\\
\bottomrule
\end{tabular}
\end{table}

\subsubsection{P1 vs.~Revisitar (I2)}\label{p1-vs.-revisitar-i2}

\begin{table}[!h]
\centering
\caption{\label{tab:cont-p1-revis}Tabla de contingencia: P1 (Sabor) vs. Revisitar (I2)}
\centering
\begin{tabular}[t]{lrrrrrr}
\toprule
  & 1 & 2 & 3 & 4 & 5 & Sum\\
\midrule
1 & 5 & 4 & 1 & 0 & 0 & 10\\
2 & 2 & 23 & 13 & 1 & 0 & 39\\
3 & 1 & 10 & 34 & 23 & 2 & 70\\
4 & 0 & 7 & 8 & 22 & 10 & 47\\
5 & 1 & 1 & 4 & 1 & 3 & 10\\
\addlinespace
Sum & 9 & 45 & 60 & 47 & 15 & 176\\
\bottomrule
\end{tabular}
\end{table}

\subsubsection{P1 vs.~Recomendar (I3)}\label{p1-vs.-recomendar-i3}

\begin{table}[!h]
\centering
\caption{\label{tab:cont-p1-recom}Tabla de contingencia: P1 (Sabor) vs. Recomendar (I3)}
\centering
\begin{tabular}[t]{lrrr}
\toprule
  & No recomendaría & Sí recomendaría & Sum\\
\midrule
1 & 10 & 0 & 10\\
2 & 33 & 6 & 39\\
3 & 30 & 40 & 70\\
4 & 12 & 35 & 47\\
5 & 2 & 8 & 10\\
\addlinespace
Sum & 87 & 89 & 176\\
\bottomrule
\end{tabular}
\end{table}

\subsection{P2 (Presentación) vs.~Variables
dependientes}\label{p2-presentaciuxf3n-vs.-variables-dependientes}

\subsubsection{P2 vs.~Satisfacción
(I1)}\label{p2-vs.-satisfacciuxf3n-i1}

\begin{table}[!h]
\centering
\caption{\label{tab:cont-p2-satisf}Tabla de contingencia: P2 (Presentación) vs. Satisfacción (I1)}
\centering
\begin{tabular}[t]{lrrrrrr}
\toprule
  & 1 & 2 & 3 & 4 & 5 & Sum\\
\midrule
1 & 0 & 4 & 0 & 0 & 0 & 4\\
2 & 4 & 8 & 9 & 1 & 0 & 22\\
3 & 1 & 18 & 24 & 25 & 1 & 69\\
4 & 1 & 6 & 11 & 36 & 4 & 58\\
5 & 0 & 3 & 3 & 12 & 5 & 23\\
\addlinespace
Sum & 6 & 39 & 47 & 74 & 10 & 176\\
\bottomrule
\end{tabular}
\end{table}

\subsubsection{P2 vs.~Revisitar (I2)}\label{p2-vs.-revisitar-i2}

\begin{table}[!h]
\centering
\caption{\label{tab:cont-p2-revis}Tabla de contingencia: P2 (Presentación) vs. Revisitar (I2)}
\centering
\begin{tabular}[t]{lrrrrrr}
\toprule
  & 1 & 2 & 3 & 4 & 5 & Sum\\
\midrule
1 & 2 & 2 & 0 & 0 & 0 & 4\\
2 & 4 & 13 & 5 & 0 & 0 & 22\\
3 & 1 & 23 & 26 & 18 & 1 & 69\\
4 & 1 & 6 & 22 & 21 & 8 & 58\\
5 & 1 & 1 & 7 & 8 & 6 & 23\\
\addlinespace
Sum & 9 & 45 & 60 & 47 & 15 & 176\\
\bottomrule
\end{tabular}
\end{table}

\subsubsection{P2 vs.~Recomendar (I3)}\label{p2-vs.-recomendar-i3}

\begin{table}[!h]
\centering
\caption{\label{tab:cont-p2-recom}Tabla de contingencia: P2 (Presentación) vs. Recomendar (I3)}
\centering
\begin{tabular}[t]{lrrr}
\toprule
  & No recomendaría & Sí recomendaría & Sum\\
\midrule
1 & 4 & 0 & 4\\
2 & 20 & 2 & 22\\
3 & 39 & 30 & 69\\
4 & 19 & 39 & 58\\
5 & 5 & 18 & 23\\
\addlinespace
Sum & 87 & 89 & 176\\
\bottomrule
\end{tabular}
\end{table}

\subsection{P3 (Frescura) vs.~Variables
dependientes}\label{p3-frescura-vs.-variables-dependientes}

\subsubsection{P3 vs.~Satisfacción
(I1)}\label{p3-vs.-satisfacciuxf3n-i1}

\begin{table}[!h]
\centering
\caption{\label{tab:cont-p3-satisf}Tabla de contingencia: P3 (Frescura) vs. Satisfacción (I1)}
\centering
\begin{tabular}[t]{lrrrrrr}
\toprule
  & 1 & 2 & 3 & 4 & 5 & Sum\\
\midrule
1 & 2 & 6 & 0 & 0 & 0 & 8\\
2 & 1 & 9 & 7 & 10 & 0 & 27\\
3 & 2 & 12 & 29 & 18 & 2 & 63\\
4 & 1 & 10 & 8 & 31 & 3 & 53\\
5 & 0 & 2 & 3 & 15 & 5 & 25\\
\addlinespace
Sum & 6 & 39 & 47 & 74 & 10 & 176\\
\bottomrule
\end{tabular}
\end{table}

\subsubsection{P3 vs.~Revisitar (I2)}\label{p3-vs.-revisitar-i2}

\begin{table}[!h]
\centering
\caption{\label{tab:cont-p3-revis}Tabla de contingencia: P3 (Frescura) vs. Revisitar (I2)}
\centering
\begin{tabular}[t]{lrrrrrr}
\toprule
  & 1 & 2 & 3 & 4 & 5 & Sum\\
\midrule
1 & 4 & 2 & 2 & 0 & 0 & 8\\
2 & 1 & 11 & 10 & 4 & 1 & 27\\
3 & 1 & 19 & 24 & 16 & 3 & 63\\
4 & 2 & 12 & 17 & 17 & 5 & 53\\
5 & 1 & 1 & 7 & 10 & 6 & 25\\
\addlinespace
Sum & 9 & 45 & 60 & 47 & 15 & 176\\
\bottomrule
\end{tabular}
\end{table}

\subsubsection{P3 vs.~Recomendar (I3)}\label{p3-vs.-recomendar-i3}

\begin{table}[!h]
\centering
\caption{\label{tab:cont-p3-recom}Tabla de contingencia: P3 (Frescura) vs. Recomendar (I3)}
\centering
\begin{tabular}[t]{lrrr}
\toprule
  & No recomendaría & Sí recomendaría & Sum\\
\midrule
1 & 8 & 0 & 8\\
2 & 19 & 8 & 27\\
3 & 34 & 29 & 63\\
4 & 22 & 31 & 53\\
5 & 4 & 21 & 25\\
\addlinespace
Sum & 87 & 89 & 176\\
\bottomrule
\end{tabular}
\end{table}

\newpage

\section{Visualizaciones}\label{visualizaciones}

\subsection{Heatmap de correlaciones}\label{heatmap-de-correlaciones}

\begin{Shaded}
\begin{Highlighting}[]
\NormalTok{corr\_resultados }\SpecialCharTok{\%\textgreater{}\%}
  \FunctionTok{mutate}\NormalTok{(}
    \AttributeTok{Item =} \FunctionTok{factor}\NormalTok{(Item, }\AttributeTok{levels =} \FunctionTok{rev}\NormalTok{(}\FunctionTok{c}\NormalTok{(}\StringTok{"P1 (Sabor)"}\NormalTok{, }\StringTok{"P2 (Presentación)"}\NormalTok{, }
                                        \StringTok{"P3 (Frescura)"}\NormalTok{))),}
    \AttributeTok{VD   =} \FunctionTok{factor}\NormalTok{(VD, }\AttributeTok{levels =} \FunctionTok{c}\NormalTok{(}\StringTok{"Satisfacción"}\NormalTok{, }\StringTok{"Revisitar"}\NormalTok{, }\StringTok{"Recomendar"}\NormalTok{))}
\NormalTok{  ) }\SpecialCharTok{\%\textgreater{}\%}
  \FunctionTok{ggplot}\NormalTok{(}\FunctionTok{aes}\NormalTok{(}\AttributeTok{x =}\NormalTok{ VD, }\AttributeTok{y =}\NormalTok{ Item, }\AttributeTok{fill =}\NormalTok{ rho)) }\SpecialCharTok{+}
  \FunctionTok{geom\_tile}\NormalTok{(}\AttributeTok{color =} \StringTok{"white"}\NormalTok{, }\AttributeTok{linewidth =} \DecValTok{1}\NormalTok{) }\SpecialCharTok{+}
  \FunctionTok{geom\_text}\NormalTok{(}\FunctionTok{aes}\NormalTok{(}\AttributeTok{label =} \FunctionTok{format}\NormalTok{(rho, }\AttributeTok{nsmall =} \DecValTok{3}\NormalTok{)), }\AttributeTok{size =} \DecValTok{4}\NormalTok{, }\AttributeTok{fontface =} \StringTok{"bold"}\NormalTok{) }\SpecialCharTok{+}
  \FunctionTok{scale\_fill\_gradient2}\NormalTok{(}\AttributeTok{low =} \StringTok{"\#d73027"}\NormalTok{, }\AttributeTok{mid =} \StringTok{"\#ffffbf"}\NormalTok{, }\AttributeTok{high =} \StringTok{"\#1a9850"}\NormalTok{,}
                        \AttributeTok{midpoint =} \DecValTok{0}\NormalTok{, }\AttributeTok{limits =} \FunctionTok{c}\NormalTok{(}\SpecialCharTok{{-}}\DecValTok{1}\NormalTok{, }\DecValTok{1}\NormalTok{),}
                        \AttributeTok{name =} \FunctionTok{expression}\NormalTok{(rho[s])) }\SpecialCharTok{+}
  \FunctionTok{labs}\NormalTok{(}\AttributeTok{x =} \ConstantTok{NULL}\NormalTok{, }\AttributeTok{y =} \ConstantTok{NULL}\NormalTok{) }\SpecialCharTok{+}
  \FunctionTok{theme\_minimal}\NormalTok{(}\AttributeTok{base\_size =} \DecValTok{11}\NormalTok{) }\SpecialCharTok{+}
  \FunctionTok{theme}\NormalTok{(}\AttributeTok{panel.grid =} \FunctionTok{element\_blank}\NormalTok{(),}
        \AttributeTok{axis.text  =} \FunctionTok{element\_text}\NormalTok{(}\AttributeTok{size =} \DecValTok{10}\NormalTok{),}
        \AttributeTok{legend.position =} \StringTok{"right"}\NormalTok{)}
\end{Highlighting}
\end{Shaded}

\begin{figure}[H]

{\centering \includegraphics[width=0.95\linewidth]{03_objetivo4_factor_sensorial_files/figure-latex/heatmap-corr-1} 

}

\caption{Correlaciones de Spearman entre los ítems del Producto (P1, P2, P3) y las variables dependientes.}\label{fig:heatmap-corr}
\end{figure}

\subsection{Gráfico de mosaico: Ítems del producto
vs.~Recomendar}\label{gruxe1fico-de-mosaico-uxedtems-del-producto-vs.-recomendar}

Los gráficos de mosaico muestran la proporción de comensales que
recomendarían (o no) en cada nivel de respuesta del ítem. El ancho de
cada columna es proporcional a la frecuencia del nivel del ítem, y la
altura de cada segmento representa la proporción de Sí/No.

\begin{Shaded}
\begin{Highlighting}[]
\NormalTok{datos }\SpecialCharTok{\%\textgreater{}\%}
  \FunctionTok{count}\NormalTok{(}\AttributeTok{P1 =} \FunctionTok{factor}\NormalTok{(P1), I3\_factor) }\SpecialCharTok{\%\textgreater{}\%}
  \FunctionTok{group\_by}\NormalTok{(P1) }\SpecialCharTok{\%\textgreater{}\%}
  \FunctionTok{mutate}\NormalTok{(}\AttributeTok{pct =}\NormalTok{ n }\SpecialCharTok{/} \FunctionTok{sum}\NormalTok{(n)) }\SpecialCharTok{\%\textgreater{}\%}
  \FunctionTok{ggplot}\NormalTok{(}\FunctionTok{aes}\NormalTok{(}\AttributeTok{x =}\NormalTok{ P1, }\AttributeTok{y =}\NormalTok{ pct, }\AttributeTok{fill =}\NormalTok{ I3\_factor)) }\SpecialCharTok{+}
  \FunctionTok{geom\_col}\NormalTok{(}\AttributeTok{width =} \FloatTok{0.7}\NormalTok{, }\AttributeTok{alpha =} \FloatTok{0.85}\NormalTok{) }\SpecialCharTok{+}
  \FunctionTok{geom\_text}\NormalTok{(}\FunctionTok{aes}\NormalTok{(}\AttributeTok{label =} \FunctionTok{paste0}\NormalTok{(}\FunctionTok{round}\NormalTok{(pct }\SpecialCharTok{*} \DecValTok{100}\NormalTok{, }\DecValTok{0}\NormalTok{), }\StringTok{"\%"}\NormalTok{)),}
            \AttributeTok{position =} \FunctionTok{position\_stack}\NormalTok{(}\AttributeTok{vjust =} \FloatTok{0.5}\NormalTok{), }\AttributeTok{size =} \DecValTok{3}\NormalTok{) }\SpecialCharTok{+}
  \FunctionTok{scale\_fill\_manual}\NormalTok{(}\AttributeTok{values =} \FunctionTok{c}\NormalTok{(}\StringTok{"No recomendaría"} \OtherTok{=} \StringTok{"\#e74c3c"}\NormalTok{,}
                                \StringTok{"Sí recomendaría"} \OtherTok{=} \StringTok{"\#27ae60"}\NormalTok{),}
                    \AttributeTok{name =} \ConstantTok{NULL}\NormalTok{) }\SpecialCharTok{+}
  \FunctionTok{scale\_y\_continuous}\NormalTok{(}\AttributeTok{labels =}\NormalTok{ scales}\SpecialCharTok{::}\FunctionTok{percent\_format}\NormalTok{()) }\SpecialCharTok{+}
  \FunctionTok{labs}\NormalTok{(}\AttributeTok{x =} \StringTok{"P1 (Sabor)"}\NormalTok{, }\AttributeTok{y =} \StringTok{"Proporción"}\NormalTok{) }\SpecialCharTok{+}
  \FunctionTok{theme\_minimal}\NormalTok{(}\AttributeTok{base\_size =} \DecValTok{10}\NormalTok{) }\SpecialCharTok{+}
  \FunctionTok{theme}\NormalTok{(}\AttributeTok{legend.position =} \StringTok{"bottom"}\NormalTok{)}
\end{Highlighting}
\end{Shaded}

\begin{figure}[H]

{\centering \includegraphics[width=0.95\linewidth]{03_objetivo4_factor_sensorial_files/figure-latex/mosaico-p1-1} 

}

\caption{Proporción de Recomendar por nivel de P1 (Sabor).}\label{fig:mosaico-p1}
\end{figure}

\begin{Shaded}
\begin{Highlighting}[]
\NormalTok{datos }\SpecialCharTok{\%\textgreater{}\%}
  \FunctionTok{count}\NormalTok{(}\AttributeTok{P2 =} \FunctionTok{factor}\NormalTok{(P2), I3\_factor) }\SpecialCharTok{\%\textgreater{}\%}
  \FunctionTok{group\_by}\NormalTok{(P2) }\SpecialCharTok{\%\textgreater{}\%}
  \FunctionTok{mutate}\NormalTok{(}\AttributeTok{pct =}\NormalTok{ n }\SpecialCharTok{/} \FunctionTok{sum}\NormalTok{(n)) }\SpecialCharTok{\%\textgreater{}\%}
  \FunctionTok{ggplot}\NormalTok{(}\FunctionTok{aes}\NormalTok{(}\AttributeTok{x =}\NormalTok{ P2, }\AttributeTok{y =}\NormalTok{ pct, }\AttributeTok{fill =}\NormalTok{ I3\_factor)) }\SpecialCharTok{+}
  \FunctionTok{geom\_col}\NormalTok{(}\AttributeTok{width =} \FloatTok{0.7}\NormalTok{, }\AttributeTok{alpha =} \FloatTok{0.85}\NormalTok{) }\SpecialCharTok{+}
  \FunctionTok{geom\_text}\NormalTok{(}\FunctionTok{aes}\NormalTok{(}\AttributeTok{label =} \FunctionTok{paste0}\NormalTok{(}\FunctionTok{round}\NormalTok{(pct }\SpecialCharTok{*} \DecValTok{100}\NormalTok{, }\DecValTok{0}\NormalTok{), }\StringTok{"\%"}\NormalTok{)),}
            \AttributeTok{position =} \FunctionTok{position\_stack}\NormalTok{(}\AttributeTok{vjust =} \FloatTok{0.5}\NormalTok{), }\AttributeTok{size =} \DecValTok{3}\NormalTok{) }\SpecialCharTok{+}
  \FunctionTok{scale\_fill\_manual}\NormalTok{(}\AttributeTok{values =} \FunctionTok{c}\NormalTok{(}\StringTok{"No recomendaría"} \OtherTok{=} \StringTok{"\#e74c3c"}\NormalTok{,}
                                \StringTok{"Sí recomendaría"} \OtherTok{=} \StringTok{"\#27ae60"}\NormalTok{),}
                    \AttributeTok{name =} \ConstantTok{NULL}\NormalTok{) }\SpecialCharTok{+}
  \FunctionTok{scale\_y\_continuous}\NormalTok{(}\AttributeTok{labels =}\NormalTok{ scales}\SpecialCharTok{::}\FunctionTok{percent\_format}\NormalTok{()) }\SpecialCharTok{+}
  \FunctionTok{labs}\NormalTok{(}\AttributeTok{x =} \StringTok{"P2 (Presentación)"}\NormalTok{, }\AttributeTok{y =} \StringTok{"Proporción"}\NormalTok{) }\SpecialCharTok{+}
  \FunctionTok{theme\_minimal}\NormalTok{(}\AttributeTok{base\_size =} \DecValTok{10}\NormalTok{) }\SpecialCharTok{+}
  \FunctionTok{theme}\NormalTok{(}\AttributeTok{legend.position =} \StringTok{"bottom"}\NormalTok{)}
\end{Highlighting}
\end{Shaded}

\begin{figure}[H]

{\centering \includegraphics[width=0.95\linewidth]{03_objetivo4_factor_sensorial_files/figure-latex/mosaico-p2-1} 

}

\caption{Proporción de Recomendar por nivel de P2 (Presentación).}\label{fig:mosaico-p2}
\end{figure}

\begin{Shaded}
\begin{Highlighting}[]
\NormalTok{datos }\SpecialCharTok{\%\textgreater{}\%}
  \FunctionTok{count}\NormalTok{(}\AttributeTok{P3 =} \FunctionTok{factor}\NormalTok{(P3), I3\_factor) }\SpecialCharTok{\%\textgreater{}\%}
  \FunctionTok{group\_by}\NormalTok{(P3) }\SpecialCharTok{\%\textgreater{}\%}
  \FunctionTok{mutate}\NormalTok{(}\AttributeTok{pct =}\NormalTok{ n }\SpecialCharTok{/} \FunctionTok{sum}\NormalTok{(n)) }\SpecialCharTok{\%\textgreater{}\%}
  \FunctionTok{ggplot}\NormalTok{(}\FunctionTok{aes}\NormalTok{(}\AttributeTok{x =}\NormalTok{ P3, }\AttributeTok{y =}\NormalTok{ pct, }\AttributeTok{fill =}\NormalTok{ I3\_factor)) }\SpecialCharTok{+}
  \FunctionTok{geom\_col}\NormalTok{(}\AttributeTok{width =} \FloatTok{0.7}\NormalTok{, }\AttributeTok{alpha =} \FloatTok{0.85}\NormalTok{) }\SpecialCharTok{+}
  \FunctionTok{geom\_text}\NormalTok{(}\FunctionTok{aes}\NormalTok{(}\AttributeTok{label =} \FunctionTok{paste0}\NormalTok{(}\FunctionTok{round}\NormalTok{(pct }\SpecialCharTok{*} \DecValTok{100}\NormalTok{, }\DecValTok{0}\NormalTok{), }\StringTok{"\%"}\NormalTok{)),}
            \AttributeTok{position =} \FunctionTok{position\_stack}\NormalTok{(}\AttributeTok{vjust =} \FloatTok{0.5}\NormalTok{), }\AttributeTok{size =} \DecValTok{3}\NormalTok{) }\SpecialCharTok{+}
  \FunctionTok{scale\_fill\_manual}\NormalTok{(}\AttributeTok{values =} \FunctionTok{c}\NormalTok{(}\StringTok{"No recomendaría"} \OtherTok{=} \StringTok{"\#e74c3c"}\NormalTok{,}
                                \StringTok{"Sí recomendaría"} \OtherTok{=} \StringTok{"\#27ae60"}\NormalTok{),}
                    \AttributeTok{name =} \ConstantTok{NULL}\NormalTok{) }\SpecialCharTok{+}
  \FunctionTok{scale\_y\_continuous}\NormalTok{(}\AttributeTok{labels =}\NormalTok{ scales}\SpecialCharTok{::}\FunctionTok{percent\_format}\NormalTok{()) }\SpecialCharTok{+}
  \FunctionTok{labs}\NormalTok{(}\AttributeTok{x =} \StringTok{"P3 (Frescura)"}\NormalTok{, }\AttributeTok{y =} \StringTok{"Proporción"}\NormalTok{) }\SpecialCharTok{+}
  \FunctionTok{theme\_minimal}\NormalTok{(}\AttributeTok{base\_size =} \DecValTok{10}\NormalTok{) }\SpecialCharTok{+}
  \FunctionTok{theme}\NormalTok{(}\AttributeTok{legend.position =} \StringTok{"bottom"}\NormalTok{)}
\end{Highlighting}
\end{Shaded}

\begin{figure}[H]

{\centering \includegraphics[width=0.95\linewidth]{03_objetivo4_factor_sensorial_files/figure-latex/mosaico-p3-1} 

}

\caption{Proporción de Recomendar por nivel de P3 (Frescura).}\label{fig:mosaico-p3}
\end{figure}

\subsection{Gráfico de dispersión con jitter: Ítems del producto vs.~VD
continuas}\label{gruxe1fico-de-dispersiuxf3n-con-jitter-uxedtems-del-producto-vs.-vd-continuas}

Para las variables dependientes Likert (Satisfacción y Revisitar), se
usa un gráfico de dispersión con jitter (desplazamiento aleatorio) y
línea de tendencia, dado que ambas variables son discretas.

\begin{Shaded}
\begin{Highlighting}[]
\NormalTok{datos }\SpecialCharTok{\%\textgreater{}\%}
  \FunctionTok{select}\NormalTok{(P1, P2, P3, idx\_satisfaccion) }\SpecialCharTok{\%\textgreater{}\%}
  \FunctionTok{pivot\_longer}\NormalTok{(P1}\SpecialCharTok{:}\NormalTok{P3, }\AttributeTok{names\_to =} \StringTok{"Item"}\NormalTok{, }\AttributeTok{values\_to =} \StringTok{"Puntaje"}\NormalTok{) }\SpecialCharTok{\%\textgreater{}\%}
  \FunctionTok{mutate}\NormalTok{(}\AttributeTok{Item =} \FunctionTok{case\_match}\NormalTok{(Item,}
                            \StringTok{"P1"} \SpecialCharTok{\textasciitilde{}} \StringTok{"P1 (Sabor)"}\NormalTok{,}
                            \StringTok{"P2"} \SpecialCharTok{\textasciitilde{}} \StringTok{"P2 (Presentación)"}\NormalTok{,}
                            \StringTok{"P3"} \SpecialCharTok{\textasciitilde{}} \StringTok{"P3 (Frescura)"}\NormalTok{)) }\SpecialCharTok{\%\textgreater{}\%}
  \FunctionTok{ggplot}\NormalTok{(}\FunctionTok{aes}\NormalTok{(}\AttributeTok{x =}\NormalTok{ Puntaje, }\AttributeTok{y =}\NormalTok{ idx\_satisfaccion)) }\SpecialCharTok{+}
  \FunctionTok{geom\_jitter}\NormalTok{(}\AttributeTok{width =} \FloatTok{0.2}\NormalTok{, }\AttributeTok{height =} \FloatTok{0.2}\NormalTok{, }\AttributeTok{alpha =} \FloatTok{0.4}\NormalTok{, }\AttributeTok{color =} \StringTok{"\#2c7fb8"}\NormalTok{, }\AttributeTok{size =} \FloatTok{1.5}\NormalTok{) }\SpecialCharTok{+}
  \FunctionTok{geom\_smooth}\NormalTok{(}\AttributeTok{method =} \StringTok{"lm"}\NormalTok{, }\AttributeTok{se =} \ConstantTok{TRUE}\NormalTok{, }\AttributeTok{color =} \StringTok{"\#e74c3c"}\NormalTok{, }\AttributeTok{linewidth =} \FloatTok{0.7}\NormalTok{) }\SpecialCharTok{+}
  \FunctionTok{facet\_wrap}\NormalTok{(}\SpecialCharTok{\textasciitilde{}}\NormalTok{Item) }\SpecialCharTok{+}
  \FunctionTok{labs}\NormalTok{(}\AttributeTok{x =} \StringTok{"Puntaje del ítem"}\NormalTok{, }\AttributeTok{y =} \StringTok{"Satisfacción (I1)"}\NormalTok{) }\SpecialCharTok{+}
  \FunctionTok{theme\_minimal}\NormalTok{(}\AttributeTok{base\_size =} \DecValTok{10}\NormalTok{)}
\end{Highlighting}
\end{Shaded}

\begin{figure}[H]

{\centering \includegraphics[width=0.95\linewidth]{03_objetivo4_factor_sensorial_files/figure-latex/jitter-satisf-1} 

}

\caption{Relación entre cada ítem del Producto y la Satisfacción (I1). La línea muestra la tendencia suavizada.}\label{fig:jitter-satisf}
\end{figure}

\begin{Shaded}
\begin{Highlighting}[]
\NormalTok{datos }\SpecialCharTok{\%\textgreater{}\%}
  \FunctionTok{select}\NormalTok{(P1, P2, P3, idx\_revisitar) }\SpecialCharTok{\%\textgreater{}\%}
  \FunctionTok{pivot\_longer}\NormalTok{(P1}\SpecialCharTok{:}\NormalTok{P3, }\AttributeTok{names\_to =} \StringTok{"Item"}\NormalTok{, }\AttributeTok{values\_to =} \StringTok{"Puntaje"}\NormalTok{) }\SpecialCharTok{\%\textgreater{}\%}
  \FunctionTok{mutate}\NormalTok{(}\AttributeTok{Item =} \FunctionTok{case\_match}\NormalTok{(Item,}
                            \StringTok{"P1"} \SpecialCharTok{\textasciitilde{}} \StringTok{"P1 (Sabor)"}\NormalTok{,}
                            \StringTok{"P2"} \SpecialCharTok{\textasciitilde{}} \StringTok{"P2 (Presentación)"}\NormalTok{,}
                            \StringTok{"P3"} \SpecialCharTok{\textasciitilde{}} \StringTok{"P3 (Frescura)"}\NormalTok{)) }\SpecialCharTok{\%\textgreater{}\%}
  \FunctionTok{ggplot}\NormalTok{(}\FunctionTok{aes}\NormalTok{(}\AttributeTok{x =}\NormalTok{ Puntaje, }\AttributeTok{y =}\NormalTok{ idx\_revisitar)) }\SpecialCharTok{+}
  \FunctionTok{geom\_jitter}\NormalTok{(}\AttributeTok{width =} \FloatTok{0.2}\NormalTok{, }\AttributeTok{height =} \FloatTok{0.2}\NormalTok{, }\AttributeTok{alpha =} \FloatTok{0.4}\NormalTok{, }\AttributeTok{color =} \StringTok{"\#2c7fb8"}\NormalTok{, }\AttributeTok{size =} \FloatTok{1.5}\NormalTok{) }\SpecialCharTok{+}
  \FunctionTok{geom\_smooth}\NormalTok{(}\AttributeTok{method =} \StringTok{"lm"}\NormalTok{, }\AttributeTok{se =} \ConstantTok{TRUE}\NormalTok{, }\AttributeTok{color =} \StringTok{"\#e74c3c"}\NormalTok{, }\AttributeTok{linewidth =} \FloatTok{0.7}\NormalTok{) }\SpecialCharTok{+}
  \FunctionTok{facet\_wrap}\NormalTok{(}\SpecialCharTok{\textasciitilde{}}\NormalTok{Item) }\SpecialCharTok{+}
  \FunctionTok{labs}\NormalTok{(}\AttributeTok{x =} \StringTok{"Puntaje del ítem"}\NormalTok{, }\AttributeTok{y =} \StringTok{"Revisitar (I2)"}\NormalTok{) }\SpecialCharTok{+}
  \FunctionTok{theme\_minimal}\NormalTok{(}\AttributeTok{base\_size =} \DecValTok{10}\NormalTok{)}
\end{Highlighting}
\end{Shaded}

\begin{figure}[H]

{\centering \includegraphics[width=0.95\linewidth]{03_objetivo4_factor_sensorial_files/figure-latex/jitter-revisitar-1} 

}

\caption{Relación entre cada ítem del Producto y la Intención de revisitar (I2). La línea muestra la tendencia suavizada.}\label{fig:jitter-revisitar}
\end{figure}

\newpage

\section{Resumen del Objetivo 4}\label{resumen-del-objetivo-4}

\begin{table}[!h]
\centering
\caption{\label{tab:resumen-obj4-print}Ranking de influencia de los ítems del Producto sobre las VD}
\centering
\begin{tabular}[t]{clc}
\toprule
Ranking & Ítem Producto & $|\rho_s|$ promedio\\
\midrule
1 & P1 (Sabor) & 0.540\\
2 & P2 (Presentación) & 0.478\\
3 & P3 (Frescura) & 0.376\\
\bottomrule
\end{tabular}
\end{table}

El factor sensorial más influyente en la toma de decisiones del comensal
es \textbf{P1 (Sabor)}, con un coeficiente de correlación de Spearman
promedio (en valor absoluto) de \(|\rho_s| =\) 0.54. Esto se sustenta
tanto en las correlaciones bivariadas como en los patrones observados en
las tablas de contingencia y los gráficos, donde se aprecia una relación
monotónica más pronunciada entre este ítem y las tres variables
dependientes.

\section{Exportación de gráficos y
tablas}\label{exportaciuxf3n-de-gruxe1ficos-y-tablas}

\begin{verbatim}
#> pdf 
#>   2
\end{verbatim}

\begin{verbatim}
#> pdf 
#>   2
\end{verbatim}

\end{document}
